% GENERATED FILE. Edit canonical lessons, then run npm run generate:books.
% canonical-source-hash: fnv1a64:23765f2653bc41d7
% canonical-lessons: HI-C450-jharnaa, HI-C450-ukharna, HI-C450-jotna, HI-C450-upajna, HI-C450-duhna

\chapter{On the Farm}
\label{ch:hi-on-the-farm}
\hlchaptermodality{450}

\begin{chapteropening}
\textbf{By the end of this chapter:} \emph{I can say leaves fell and a tree was uprooted, a farmer ploughed, the crop grew and the cow was milked.}

Complete the last lesson of chapter 450: I can say leaves fell and a tree was uprooted, a farmer ploughed, the crop grew and the cow was milked.
\end{chapteropening}

\section[\texorpdfstring{\hi{झड़ना} (jhaṛnā)}{झड़ना (jhaṛnā)}]{\hi{झड़ना} (jhaṛnā) — to fall off; to be shed (of leaves)}

\label{lesson:HI-C450-jharnaa}



\begin{quote}
Before the new one: say the Hindi for to graze, then the Hindi for to take to graze.
\end{quote}

\subsection*{You'll want to know: \hi{झड़ना}}
\textbf{\hi{झड़ना}} — \emph{jhaṛnā} — \textquotedblleft{}to fall off; to be shed (of leaves)\textquotedblright{}.

\begin{quote}
\hi{पतझड़} \hi{में} \hi{पेड़ों} \hi{से} \hi{पत्ते} \hi{झड़ते} \hi{हैं।}
\end{quote}

\begin{quote}
\emph{patjhaṛ meṁ peṛoṁ se patte jhaṛte haiṁ.}
\end{quote}

\begin{quote}
In autumn the leaves fall from the trees.
\end{quote}

Past (it …): \hi{वह} \hi{पत्ता} \hi{पेड़} \hi{से} \hi{झड़} \hi{गया}, \emph{vah pattā peṛ se jhaṛ gayā}.

\begin{cousinweb}
\textbf{\hi{झड़ना}} (\emph{jhaṛnā}) goes back to \emph{*jhaṭati}, \textquotedblleft{}falls\textquotedblright{}, a form CDIAL reconstructs (CDIAL 5328), which lists Hindi \emph{jhaṛnā}, \textquotedblleft{}to fall\textquotedblright{}; \hi{पतझड़}, autumn, is the leaf-fall.
\end{cousinweb}

\begin{grammarlens}[title={What you've built}]
One more word for getting around.
\end{grammarlens}

\subsection*{Your turn}
\begin{itemize}
\raggedright
  \item \emph{Say it:} \emph{jhaṛnā}
  \item \emph{Say it:} \emph{jhaṛnā}, once more
  \item \emph{Say it:} \emph{patjhaṛ meṁ peṛoṁ se patte jhaṛte haiṁ.} — In autumn the leaves fall from the trees.
  \item \emph{Recall:} say the Hindi for to graze, then the Hindi for to take to graze, then say \emph{jhaṛnā} again
\end{itemize}

\begin{tcolorbox}[breakable,colback=teal!4,colframe=teal!35!black,title={Before you move on}]
What does \hi{झड़ना} mean? (\textquotedblleft{}To fall off; to be shed (of leaves)\textquotedblright{}.) Say it once more.
\end{tcolorbox}
\section[\texorpdfstring{\hi{उखड़ना} (ukhaṛnā)}{उखड़ना (ukhaṛnā)}]{\hi{उखड़ना} (ukhaṛnā) — to be uprooted; to come loose}

\label{lesson:HI-C450-ukharna}



\begin{quote}
Before the new one: say the Hindi for to take to graze, then the Hindi for to be shed.
\end{quote}

\subsection*{You'll want to know: \hi{उखड़ना}}
\textbf{\hi{उखड़ना}} — \emph{ukhaṛnā} — \textquotedblleft{}to be uprooted; to come loose\textquotedblright{}.

\begin{quote}
\hi{तूफ़ान} \hi{में} \hi{कई} \hi{पेड़} \hi{उखड़} \hi{गए।}
\end{quote}

\begin{quote}
\emph{tūfān meṁ kaī peṛ ukhaṛ gae.}
\end{quote}

\begin{quote}
Several trees were uprooted in the storm.
\end{quote}

Past (it …): \hi{वह} \hi{खंभा} \hi{उखड़} \hi{गया}, \emph{vah khambhā ukhaṛ gayā}.

\begin{cousinweb}
\textbf{\hi{उखड़ना}} (\emph{ukhaṛnā}) goes back to \emph{*utskarati}, \textquotedblleft{}digs out\textquotedblright{}, a form CDIAL reconstructs (CDIAL 1895), which lists Hindi \emph{ukhaṛnā}, \textquotedblleft{}to be rooted up\textquotedblright{}.
\end{cousinweb}

\begin{grammarlens}[title={What you've built}]
One more word for getting around.
\end{grammarlens}

\subsection*{Your turn}
\begin{itemize}
\raggedright
  \item \emph{Say it:} \emph{ukhaṛnā}
  \item \emph{Say it:} \emph{ukhaṛnā}, once more
  \item \emph{Say it:} \emph{tūfān meṁ kaī peṛ ukhaṛ gae.} — Several trees were uprooted in the storm.
  \item \emph{Recall:} say the Hindi for to take to graze, then the Hindi for to be shed, then say \emph{ukhaṛnā} again
\end{itemize}

\begin{tcolorbox}[breakable,colback=teal!4,colframe=teal!35!black,title={Before you move on}]
What does \hi{उखड़ना} mean? (\textquotedblleft{}To be uprooted; to come loose\textquotedblright{}.) Say it once more.
\end{tcolorbox}
\section[\texorpdfstring{\hi{जोतना} (jotnā)}{जोतना (jotnā)}]{\hi{जोतना} (jotnā) — to plough}

\label{lesson:HI-C450-jotna}



\begin{quote}
Before the new one: say the Hindi for to be shed, then the Hindi for to be uprooted.
\end{quote}

\subsection*{You'll want to know: \hi{जोतना}}
\textbf{\hi{जोतना}} — \emph{jotnā} — \textquotedblleft{}to plough\textquotedblright{}.

\begin{quote}
\hi{किसान} \hi{सुबह} \hi{से} \hi{खेत} \hi{जोत} \hi{रहा} \hi{है।}
\end{quote}

\begin{quote}
\emph{kisān subah se khet jot rahā hai.}
\end{quote}

\begin{quote}
The farmer has been ploughing the field since morning.
\end{quote}

Past, with \hi{ने} (he or she …): \hi{उसने} \hi{खेत} \hi{जोता}, \emph{usne khet jotā}.

\begin{cousinweb}
\textbf{\hi{जोतना}} (\emph{jotnā}) is a tadbhava of Sanskrit \emph{yoktrayati}, \textquotedblleft{}fastens, ties\textquotedblright{} (CDIAL 10524), which lists Hindi \emph{jotnā}, \textquotedblleft{}to yoke, plough\textquotedblright{}: to yoke the oxen, and so to plough.
\end{cousinweb}

\begin{grammarlens}[title={What you've built}]
One more word for getting around.
\end{grammarlens}

\subsection*{Your turn}
\begin{itemize}
\raggedright
  \item \emph{Say it:} \emph{jotnā}
  \item \emph{Say it:} \emph{jotnā}, once more
  \item \emph{Say it:} \emph{kisān subah se khet jot rahā hai.} — The farmer has been ploughing the field since morning.
  \item \emph{Recall:} say the Hindi for to be shed, then the Hindi for to be uprooted, then say \emph{jotnā} again
\end{itemize}

\begin{tcolorbox}[breakable,colback=teal!4,colframe=teal!35!black,title={Before you move on}]
What does \hi{जोतना} mean? (\textquotedblleft{}To plough\textquotedblright{}.) Say it once more.
\end{tcolorbox}
\section[\texorpdfstring{\hi{उपजना} (upajnā)}{उपजना (upajnā)}]{\hi{उपजना} (upajnā) — to grow (of crops); to arise}

\label{lesson:HI-C450-upajna}



\begin{quote}
Before the new one: say the Hindi for to be uprooted, then the Hindi for to plough.
\end{quote}

\subsection*{You'll want to know: \hi{उपजना}}
\textbf{\hi{उपजना}} — \emph{upajnā} — \textquotedblleft{}to grow (of crops); to arise\textquotedblright{}.

\begin{quote}
\hi{इस} \hi{खेत} \hi{में} \hi{अच्छा} \hi{अनाज} \hi{उपजता} \hi{है।}
\end{quote}

\begin{quote}
\emph{is khet meṁ acchā anāj upajtā hai.}
\end{quote}

\begin{quote}
Good grain grows in this field.
\end{quote}

Past (it …): \hi{वह} \hi{फ़सल} \hi{अच्छी} \hi{उपजी}, \emph{vah fasal acchī upjī}.

\begin{cousinweb}
\textbf{\hi{उपजना}} (\emph{upajnā}) is a tadbhava of Sanskrit \emph{utpadyate}, \textquotedblleft{}arises, originates\textquotedblright{} (CDIAL 1814), which lists Hindi \emph{upajnā}, \textquotedblleft{}to originate, grow\textquotedblright{}; \hi{उपज}, a crop, is its noun.
\end{cousinweb}

\begin{grammarlens}[title={What you've built}]
One more word for getting around.
\end{grammarlens}

\subsection*{Your turn}
\begin{itemize}
\raggedright
  \item \emph{Say it:} \emph{upajnā}
  \item \emph{Say it:} \emph{upajnā}, once more
  \item \emph{Say it:} \emph{is khet meṁ acchā anāj upajtā hai.} — Good grain grows in this field.
  \item \emph{Recall:} say the Hindi for to be uprooted, then the Hindi for to plough, then say \emph{upajnā} again
\end{itemize}

\begin{tcolorbox}[breakable,colback=teal!4,colframe=teal!35!black,title={Before you move on}]
What does \hi{उपजना} mean? (\textquotedblleft{}To grow (of crops); to arise\textquotedblright{}.) Say it once more.
\end{tcolorbox}
\section[\texorpdfstring{\hi{दुहना} (duhnā)}{दुहना (duhnā)}]{\hi{दुहना} (duhnā) — to milk}

\label{lesson:HI-C450-duhna}



\begin{quote}
Before the new one: say the Hindi for to plough, then the Hindi for to grow (of crops).
\end{quote}

\subsection*{You'll want to know: \hi{दुहना}}
\textbf{\hi{दुहना}} — \emph{duhnā} — \textquotedblleft{}to milk\textquotedblright{}.

\begin{quote}
\hi{दादी} \hi{रोज़} \hi{सुबह} \hi{गाय} \hi{दुहती} \hi{हैं।}
\end{quote}

\begin{quote}
\emph{dādī roz subah gāy duhtī haiṁ.}
\end{quote}

\begin{quote}
Grandmother milks the cow every morning.
\end{quote}

Past, with \hi{ने} (he or she …): \hi{उसने} \hi{भैंस} \hi{दुही}, \emph{usne bhaĩs duhī}.

\begin{cousinweb}
\textbf{\hi{दुहना}} (\emph{duhnā}) goes back to \emph{*duhati}, \textquotedblleft{}milks\textquotedblright{}, a form CDIAL reconstructs (CDIAL 6476), which lists Hindi \emph{duhnā}, \textquotedblleft{}to milk\textquotedblright{}.
\end{cousinweb}

\begin{grammarlens}[title={What you've built}]
One more word for getting around.
\end{grammarlens}

\subsection*{Your turn}
\begin{itemize}
\raggedright
  \item \emph{Say it:} \emph{duhnā}
  \item \emph{Say it:} \emph{duhnā}, once more
  \item \emph{Say it:} \emph{dādī roz subah gāy duhtī haiṁ.} — Grandmother milks the cow every morning.
  \item \emph{Recall:} say the Hindi for to plough, then the Hindi for to grow (of crops), then say \emph{duhnā} again
\end{itemize}

\begin{tcolorbox}[breakable,colback=teal!4,colframe=teal!35!black,title={Before you move on}]
What does \hi{दुहना} mean? (\textquotedblleft{}To milk\textquotedblright{}.) Say it once more.
\end{tcolorbox}
